% ========================================================
% BÀI TẬP TOÁN 5 - KẾT NỐI TRI THỨC VỚI CUỘC SỐNG
% Bài 1: Ôn tập số tự nhiên
% Định dạng: BTTeX Part III - Trắc nghiệm điền số ngắn
% ========================================================

% Câu 1
\begin{ex}
  Cho số tự nhiên $85\,234\,761$. Chữ số thuộc hàng chục triệu của số này là bao nhiêu?
  \shortans{8}
  \loigiai{
    Số $85\,234\,761$ có lớp triệu gồm hai chữ số: chữ số 8 thuộc hàng chục triệu và chữ số 5 thuộc hàng triệu. Vậy chữ số hàng chục triệu là 8.
  }
\end{ex}

% Câu 2
\begin{ex}
  Trong số $7\,492\,315$, chữ số $4$ có giá trị bằng bao nhiêu?
  \shortans{400000}
  \loigiai{
    Trong số $7\,492\,315$, chữ số 4 nằm ở hàng trăm nghìn nên có giá trị là $400\,000$.
  }
\end{ex}

% Câu 3
\begin{ex}
  Số tự nhiên gồm $6$ triệu, $5$ trăm nghìn, $8$ nghìn và $4$ đơn vị được viết là bao nhiêu?
  \shortans{6508004}
  \loigiai{
    Số gồm 6 triệu, 5 trăm nghìn, 0 chục nghìn, 8 nghìn, 0 trăm, 0 chục và 4 đơn vị được viết là $6\,508\,004$.
  }
\end{ex}

% Câu 4
\begin{ex}
  Tìm số tự nhiên thích hợp điền vào chỗ chấm:
  $3\,450\,720 = 3\,000\,000 + 400\,000 + \dots + 700 + 20$.
  \shortans{50000}
  \loigiai{
    Phân tích số $3\,450\,720$ theo các hàng ta có chữ số 5 thuộc hàng chục nghìn, có giá trị là $50\,000$. Do đó số cần điền vào chỗ chấm là $50\,000$.
  }
\end{ex}

% Câu 5
\begin{ex}
  Tìm số tự nhiên liền trước của số lớn nhất có sáu chữ số.
  \shortans{999998}
  \loigiai{
    Số lớn nhất có sáu chữ số là $999\,999$. Số liền trước của nó là $999\,999 - 1 = 999\,998$.
  }
\end{ex}

% Câu 6
\begin{ex}
  Số tự nhiên lẻ bé nhất có năm chữ số khác nhau là bao nhiêu?
  \shortans{10235}
  \loigiai{
    Để có số bé nhất có năm chữ số khác nhau: chữ số hàng chục nghìn là 1, hàng nghìn là 0, hàng trăm là 2, hàng chục là 3. Để số là số lẻ, chữ số hàng đơn vị phải là số lẻ nhỏ nhất khác các chữ số trên, đó là chữ số 5. Vậy số cần tìm là $10\,235$.
  }
\end{ex}

% Câu 7
\begin{ex}
  Từ $100$ đến $999$ có tất cả bao nhiêu số tự nhiên chẵn?
  \shortans{450}
  \loigiai{
    Các số chẵn từ 100 đến 999 tạo thành dãy cách đều 2 đơn vị: $100, 102, 104, \dots, 998$.
    Số lượng số chẵn là: $(998 - 100) : 2 + 1 = 450$ (số).
  }
\end{ex}

% Câu 8
\begin{ex}
  Tìm chữ số $x$ lớn nhất để số có sáu chữ số thỏa mãn: $\overline{45x\,820} < 456\,820$.
  \shortans{5}
  \loigiai{
    Hai số có cùng số chữ số, chữ số hàng trăm nghìn (4) và hàng chục nghìn (5) bằng nhau. Để $\overline{45x\,820} < 456\,820$ thì chữ số hàng nghìn phải thỏa mãn $x < 6$. Chữ số lớn nhất thỏa mãn là 5.
  }
\end{ex}

% Câu 9
\begin{ex}
  Làm tròn số $48\,725$ đến hàng nghìn ta được số nào?
  \shortans{49000}
  \loigiai{
    Khi làm tròn số $48\,725$ đến hàng nghìn, ta xét chữ số hàng trăm là 7. Vì $7 \ge 5$ nên ta làm tròn lên, được số $49\,000$.
  }
\end{ex}

% Câu 10
\begin{ex}
  Làm tròn số $2\,849\,150$ đến hàng trăm nghìn ta được số nào?
  \shortans{2800000}
  \loigiai{
    Khi làm tròn số $2\,849\,150$ đến hàng trăm nghìn, ta xét chữ số hàng chục nghìn là 4. Vì $4 < 5$ nên ta làm tròn xuống, được số $2\,800\,000$.
  }
\end{ex}

% Câu 11
\begin{ex}
  Dùng cả năm chữ số $0; 2; 5; 7; 9$ (mỗi chữ số dùng đúng một lần) để viết số tự nhiên chẵn lớn nhất có năm chữ số. Số viết được là bao nhiêu?
  \shortans{97520}
  \loigiai{
    Để số viết được là số chẵn lớn nhất có năm chữ số: chữ số tận cùng phải là 0 hoặc 2. Nếu tận cùng là 0, ta xếp các chữ số còn lại từ lớn đến bé được số $97\,520$. Nếu tận cùng là 2, ta được số $97\,502$. Vì $97\,520 > 97\,502$ nên số lớn nhất là $97\,520$.
  }
\end{ex}

% Câu 12
\begin{ex}
  Từ bốn chữ số $1; 0; 4; 8$, viết được số lớn nhất có bốn chữ số khác nhau và số bé nhất có bốn chữ số khác nhau. Hiệu của hai số đó là bao nhiêu?
  \shortans{7362}
  \loigiai{
    Số lớn nhất có bốn chữ số khác nhau viết từ 4 chữ số trên là $8\,410$.
    Số bé nhất có bốn chữ số khác nhau viết từ 4 chữ số trên là $1\,048$.
    Hiệu của hai số đó là: $8\,410 - 1\,048 = 7\,362$.
  }
\end{ex}

% Câu 13
\begin{ex}
  Có bao nhiêu số tròn chục nghìn có sáu chữ số và nhỏ hơn $300\,000$?
  \shortans{20}
  \loigiai{
    Các số tròn chục nghìn có sáu chữ số và nhỏ hơn $300\,000$ là: $100\,000, 110\,000, 120\,000, \dots, 290\,000$.
    Số lượng số là: $(290\,000 - 100\,000) : 10\,000 + 1 = 20$ (số).
  }
\end{ex}

% Câu 14
\begin{ex}
  Trong các số tự nhiên từ $1$ đến $100$, có bao nhiêu số chia hết cho cả $2$ và $5$?
  \shortans{10}
  \loigiai{
    Số chia hết cho cả 2 và 5 phải có chữ số tận cùng là 0 (chia hết cho 10). Các số đó từ 1 đến 100 là: $10, 20, 30, \dots, 100$.
    Số lượng số là: $(100 - 10) : 10 + 1 = 10$ (số).
  }
\end{ex}

% Câu 15
\begin{ex}
  Một số tự nhiên có năm chữ số, biết rằng khi xóa đi chữ số $7$ ở hàng chục nghìn của số đó thì số đó giảm đi bao nhiêu đơn vị?
  \shortans{70000}
  \loigiai{
    Chữ số 7 ở hàng chục nghìn có giá trị bằng $70\,000$. Do đó, khi xóa chữ số 7 này thì số đó sẽ giảm đi đúng $70\,000$ đơn vị.
  }
\end{ex}

% Câu 16
\begin{ex}
  Tìm số tự nhiên bé nhất có các chữ số khác nhau mà tổng các chữ số của nó bằng $21$.
  \shortans{489}
  \loigiai{
    Để số là bé nhất, số chữ số phải ít nhất. Ta chọn các chữ số ở hàng thấp lớn nhất có thể: hàng đơn vị là 9, hàng chục là 8. Chữ số hàng trăm là: $21 - (9 + 8) = 4$. Vậy số bé nhất thỏa mãn là $489$.
  }
\end{ex}

% Câu 17
\begin{ex}
  Tìm số tự nhiên lớn nhất có bốn chữ số khác nhau mà tổng các chữ số của nó bằng $10$.
  \shortans{7210}
  \loigiai{
    Để số lớn nhất có bốn chữ số khác nhau thì chữ số hàng nghìn phải lớn nhất có thể. Chọn chữ số hàng nghìn là 7, khi đó tổng ba chữ số còn lại là $10 - 7 = 3$. Ba chữ số khác nhau có tổng bằng 3 chỉ có thể là $2; 1; 0$. Sắp xếp giảm dần ta được số lớn nhất là $7\,210$.
  }
\end{ex}

% Câu 18
\begin{ex}
  Để đánh số trang một cuốn sách dày $45$ trang (từ trang $1$ đến trang $45$), người ta cần dùng tất cả bao nhiêu chữ số?
  \shortans{81}
  \loigiai{
    Từ trang 1 đến trang 9 có: $(9 - 1) + 1 = 9$ trang có 1 chữ số, cần $9 \times 1 = 9$ chữ số.
    Từ trang 10 đến trang 45 có: $45 - 10 + 1 = 36$ trang có 2 chữ số, cần $36 \times 2 = 72$ chữ số.
    Tổng số chữ số cần dùng là: $9 + 72 = 81$ chữ số.
  }
\end{ex}

% Câu 19
\begin{ex}
  Người ta dùng $189$ chữ số để đánh số trang một cuốn sách bắt đầu từ trang $1$. Hỏi cuốn sách đó dày bao nhiêu trang?
  \shortans{99}
  \loigiai{
    Số chữ số để đánh 9 trang đầu (từ trang 1 đến trang 9) là $9 \times 1 = 9$ chữ số.
    Số chữ số còn lại là: $189 - 9 = 180$ chữ số.
    Vì các trang tiếp theo đều có 2 chữ số nên số trang có 2 chữ số là: $180 : 2 = 90$ trang.
    Vậy cuốn sách dày số trang là: $9 + 90 = 99$ trang.
  }
\end{ex}

% Câu 20
\begin{ex}
  Cho dãy số cách đều: $3, 7, 11, 15, 19, \dots$. Số thứ $20$ của dãy số này là bao nhiêu?
  \shortans{79}
  \loigiai{
    Dãy số cách đều có khoảng cách giữa hai số liên tiếp là $7 - 3 = 4$ đơn vị.
    Số thứ 20 của dãy là: $3 + (20 - 1) \times 4 = 3 + 76 = 79$.
  }
\end{ex}

% Câu 21
\begin{ex}
  Một số tự nhiên khi làm tròn đến hàng trăm nghìn thì được số $600\,000$. Số tự nhiên lớn nhất có thể thỏa mãn điều kiện đó là bao nhiêu?
  \shortans{649999}
  \loigiai{
    Các số tự nhiên khi làm tròn đến hàng trăm nghìn được $600\,000$ là các số từ $550\,000$ đến $649\,999$. Vì vậy, số tự nhiên lớn nhất có thể là $649\,999$.
  }
\end{ex}

% Câu 22
\begin{ex}
  Cho một số có ba chữ số mà chữ số hàng trăm là $4$. Nếu chuyển chữ số $4$ xuống hàng đơn vị và giữ nguyên thứ tự hai chữ số còn lại thì được số mới bé hơn số ban đầu $162$ đơn vị. Số ban đầu là bao nhiêu?
  \shortans{426}
  \loigiai{
    Gọi số ban đầu là $\overline{4ab} = 400 + \overline{ab}$.
    Khi chuyển chữ số 4 xuống hàng đơn vị, ta được số mới là $\overline{ab4} = 10 \times \overline{ab} + 4$.
    Theo đề bài: $(400 + \overline{ab}) - (10 \times \overline{ab} + 4) = 162$.
    Suy ra: $396 - 9 \times \overline{ab} = 162 \implies 9 \times \overline{ab} = 234 \implies \overline{ab} = 26$.
    Vậy số ban đầu là $426$.
  }
\end{ex}

% Câu 23
\begin{ex}
  Có bao nhiêu số tự nhiên có ba chữ số mà tổng các chữ số của nó bằng $3$?
  \shortans{6}
  \loigiai{
    Phân tích số 3 thành tổng của ba chữ số:
    - Nhóm $\{3, 0, 0\}$: lập được 1 số là $300$.
    - Nhóm $\{2, 1, 0\}$: lập được các số là $210, 201, 120, 102$ (4 số).
    - Nhóm $\{1, 1, 1\}$: lập được 1 số là $111$.
    Tổng cộng có: $1 + 4 + 1 = 6$ số.
  }
\end{ex}

% Câu 24
\begin{ex}
  Tích $1 \times 2 \times 3 \times 4 \times 5 \times 6 \times 7 \times 8 \times 9 \times 10$ có tận cùng là bao nhiêu chữ số $0$?
  \shortans{2}
  \loigiai{
    Mỗi chữ số 0 tận cùng được tạo thành bởi một tích của thừa số chia hết cho 5 với một số chẵn. Trong tích từ 1 đến 10 có hai thừa số chia hết cho 5 là 5 và 10 ($10 = 2 \times 5$). Do đó tích có đúng 2 chữ số 0 tận cùng.
  }
\end{ex}

% Câu 25
\begin{ex}
  Tìm một số tự nhiên có hai chữ số, biết rằng số đó gấp $9$ lần tổng các chữ số của nó.
  \shortans{81}
  \loigiai{
    Gọi số cần tìm là $\overline{ab}$ ($a \ne 0; a, b < 10$).
    Theo đề bài ta có:
    $\overline{ab} = 9 \times (a + b)$
    $10a + b = 9a + 9b$
    $a = 8b$.
    Vì $a$ là chữ số khác 0 nên chỉ có trường hợp $b = 1$ suy ra $a = 8$.
    Vậy số tự nhiên cần tìm là $81$.
  }
\end{ex}
